\documentclass[letterpaper, 12pt, parskip=full,DIV=10]{scrartcl}

\title{SOSC 13200: Social Science Inquiry II.}
\subtitle{University of Chicago, Winter 2027.}

\RequirePackage{assets/template_MOW}
\usepackage{xurl}
\usepackage{needspace}
\clubpenalty=10000
\widowpenalty=10000

\begin{document}
\nobibliography{assets/bib}
\maketitle

\vskip -10 ex

\textbf{Location:} Campus North Residential Commons 158\\
\textbf{Course time:} Tue/Thu 12:30--13:50 (Section 7); 14:00--15:20 (Section 9)\\
\textbf{Course github:} \href{https://github.com/UChicago-pol-methods/SOSC-13200-W27-7-9}{\texttt{UChicago-pol-methods/SOSC-13200-W27-7-9}}

\textbf{Instructor:} Molly Offer-Westort; \href{mailto:mollyow@uchicago.edu}{mollyow@uchicago.edu}\\
\textbf{Office hours:} Book at \url{https://calendar.app.google/HA29eTYFs5qgeHFT8}\\
\textbf{Office:} Pick Hall 328

\paragraph{Course description and objectives.}
In this course, you will learn to approach data thinking like a social scientist.
That means thinking about where your data comes from and how it is measured, and assessing and describing relationships among variables.
Throughout the course, you will be exposed to statistical tests and other quantitative methods used in social science research.
You will also work with statistical software for data analysis and data visualization.

This session of the course focuses on the theme of the ``credibility revolution'' in the social sciences:
a move toward empirical studies built with deliberate data collection, rigorous research design, and appropriate strategies for analysis.
The impact of this work was recognized by the 2021 Nobel prize in economics, awarded to David Card, Joshua Angrist, and Guido Imbens.
We will also consider the contributions of carefully designed field experiments, and the work of 2019 economics Nobel prize winners, Abhijit Banerjee, Esther Duflo, and Michael Kremer.

\Needspace{7\baselineskip}
\paragraph{Textbooks.}
\begin{itemize}
\item Probability \& Statistics:
\begin{itemize}
\item \bibentry{stark1997sticigui}. Chapter and section assignments appear alongside each class in the course outline; additional references are marked optional. Interactive exercises and videos are optional. Contents: \url{https://www.stat.berkeley.edu/~stark/SticiGui/Text/toc.htm}\\
This textbook is freely available online. 
\end{itemize}
\item Social science research: \bibentry{buenodemesquita2021thinking}. Chapter 5.
\item Coding in R:
\begin{itemize}
\item \bibentry{wickhamggplot2online}. Online book: \url{https://ggplot2-book.org/}
\end{itemize}
\end{itemize}

\paragraph{Mathematical proficiency.}
It is not assumed that students have had prior exposure to probability and statistics, and this course will not use calculus or linear algebra.
Familiarity with these topics may be useful in particular for some of the extra reference readings.

\paragraph{Statistical software.}
The statistical software used in this iteration of this class will be R, which we will use along with the RStudio interface.
R is a statistical language and environment for data manipulation, calculation, and visualization; the software for R and RStudio are free to download online.
R can be a very flexible and powerful tool, and it is widely used in statistical and social science research, as well as in some industry settings.
Instructors in other sessions of this course may use Stata, which is also widely used and relevant in particular in the fields of economics and public policy.
In general, basic familiarity with R, Stata, and python can be useful for research in the social sciences.
You are not assumed to have prior experience with R for this class.

\paragraph{Computing.}
You will need access to a computer throughout the quarter with R and RStudio installed.
If you do not have access to a private computer to install this software, University library computers are equipped with RStudio, and you should be able to save your assignments on Box while using these computers.
If you have any issues with access to a computer, please reach out to me and we will find a solution.

\paragraph{Course logistics.}
The course discussion platform will be announced at the start of the quarter.
Lecture materials, problem sets, and in-class code will be posted on the \href{https://github.com/UChicago-pol-methods/SOSC-13200-W27-7-9}{course GitHub page}.
Problem set solutions will be posted after the due date on Canvas.

\paragraph{Assignments.}
Homework consists of eight assignments: two R analyses with compiled reports and six written assignments.
Homeworks are graded on a 0--10 scale, with 80 raw points total.
Homeworks are worth 20\% of the course grade. Participation is worth 10\%. The in-class midterm is worth 30\%, and the cumulative in-person final is worth 40\%.
There is no take-home midterm or take-home final.

\begin{center}
\begin{tabular}{l l l}
\textbf{Component} & \textbf{Weight} & \textbf{Due date} \\
\hline
Participation & 10\% & -- \\
Homework 1 & 2.5\% & Fri, January 8 (11:59pm) \\
Homework 2 & 2.5\% & Fri, January 15 (11:59pm) \\
Homework 3 & 2.5\% & Fri, January 22 (11:59pm) \\
Homework 4 & 2.5\% & Fri, January 29 (11:59pm) \\
Homework 5 & 2.5\% & Fri, February 12 (11:59pm) \\
Homework 6 & 2.5\% & Fri, February 19 (11:59pm) \\
Homework 7 & 2.5\% & Fri, February 26 (11:59pm) \\
Homework 8 & 2.5\% & Tue, March 2 (9:00am) \\
Midterm & 30\% & Thu, February 4 (in class)\\
Final (cumulative) & 40\% & Thu, March 4 (in class)\\
\end{tabular}
\end{center}

The midterm covers Classes 1.1--4.2 and Homeworks 1--4. 
There is no homework due during midterm week.
The cumulative final covers the whole course. 
Both exams are planned for 80 minutes.
You may bring a non-programmable calculator to exams. 

Submit assignments to Canvas as one readable PDF; Homeworks 2 and 5 also require your R script submitted as a separate .R file. 
Write your name and section on each submission.
Handwritten scans are welcome if legible, except for code-based assignments. 
Unless an extension is arranged, late assignments lose one raw point per 24 hours for up to three days; no credit is given after the third day. 
Extensions due to previously scheduled conflicts must be arranged at least one week before the assignment is due. 
Extensions due to personal emergencies will be arranged on a case-by-case basis, but must be arranged before the due date. 

Participation is assessed through preparation, substantive engagement with evidence, and brief individual written exercises.
%The best fourteen of sixteen non-exam activity submissions count; speaking frequently is not required. 
Computers and phones may not be used during ordinary class activities unless requested by the instructor or needed for an accommodation. 
R work is completed outside class.

\paragraph{Accommodations.}
Please reach out to me directly if you would like to request accommodations for the course to better facilitate your learning.
Student Disability Services (\url{https://disabilities.uchicago.edu/}) is also available to provide you resources and support, and may provide approval for specific academic accommodations.
Informing me in a timely manner will help me to ensure accommodations are met and I am able to implement an appropriate assessment of your learning.

\paragraph{Academic honesty.}
All students are expected to be familiar with and follow the University\textquotesingle s academic honesty policies (\url{https://studentmanual.uchicago.edu/academic-policies/academic-honesty-plagiarism/}).
You are responsible for understanding the code and ensuring it works correctly.
On the two R assignments, AI may explain an error or help debug after your own attempt. Compose written interpretations yourself, disclose assistance, and verify the result. 

Article and book-chapter readings will be provided in the course reader. Links are included for online materials.

\section*{Course outline.}

\subsection*{Week 1: Course introduction.}
\subsubsection*{Class 1.1 (Tuesday 1/5).}
\begin{itemize}
\item Overview of course objective: how to formalize and test a theory with data and statistical tools.
\item Some motivating examples.
\item Introduction to coding: what the software will be, how you will download it, basics of setting yourself up with good directory hygiene.
\end{itemize}

\subsubsection*{Readings.}
\begin{itemize}
\item \bibentry{tabarrok2021credibility}. \href{https://marginalrevolution.com/marginalrevolution/2021/10/the-credibility-revolution-1.html}{Read online.}
\end{itemize}

\subsubsection*{Class 1.2 (Thursday 1/7).}
\begin{itemize}
\item What kinds of questions can we want to ask in social science research?
How do we start formalizing a question that is tractable and falsifiable?
What population is the question in reference to?
\end{itemize}

\subsubsection*{Readings.}
\begin{itemize}
\item \bibentry{holland1986statistics}. Sections 1--4, 6--7 and 8.1 (Causation and Medicine); skip section 5, section 8.2, and the rejoinders.
\end{itemize}

\textbf{Homework 1 due Friday 1/8 at 11:59pm:}
Counterfactuals and credible comparisons. Use Holland to analyze potential outcomes, random assignment, and the assumptions behind a natural comparison.

\subsection*{Week 2: Summarizing data numerically and visually.}
\subsubsection*{Class 2.1 (Tuesday 1/12).}
\begin{itemize}
\item Working with the Card and Krueger dataset, summarizing data numerically.
Summaries of univariate data (sample mean, sample median, quantiles, sample variance, sample standard deviation).
\end{itemize}

\subsubsection*{Readings.}
\begin{itemize}
\item \bibentry{cardkrueger1994minimum}. Introduction and data description, Tables 1--2 and the employment definition on p. 775.\begin{itemize}
\item We will come back to this paper again in Week 7 to discuss the content and results; for this first pass, we will just focus on the data.
\end{itemize}
\item \textbf{Optional:} Stark, \textit{SticiGui}, \href{https://www.stat.berkeley.edu/~stark/SticiGui/Text/histograms.htm}{Chapter 3: Statistics}: ``Data'' and ``Percentiles and Quartiles''; \href{https://www.stat.berkeley.edu/~stark/SticiGui/Text/location.htm}{Chapter 4: Measures of Location and Spread}: ``Measures of Location,'' ``Spread or Variability'' and ``The Range, IQR and SD.''
\end{itemize}

\subsubsection*{Class 2.2 (Thursday 1/14).}
\begin{itemize}
\item Working with the Card and Krueger dataset, summarizing data visually.
\item Graphical representation of different types of univariate data: histograms, density plots, boxplot, bar charts.
\item Challenges with measurement.
\end{itemize}

\subsubsection*{Readings.}
\begin{itemize}
\item \bibentry{adcockcollier2001measurement}. Introductory discussion.
\item \textbf{Optional:} Stark, \textit{SticiGui}, \href{https://www.stat.berkeley.edu/~stark/SticiGui/Text/histograms.htm}{Chapter 3: Statistics}: ``Histograms''; \href{https://www.stat.berkeley.edu/~stark/SticiGui/Text/scatterplots.htm}{Chapter 5: Multivariate Data and Scatterplots}: ``Scatterplots'' and ``Describing Scatterplots.''
\end{itemize}

\textbf{Homework 2 due Friday 1/15 at 11:59pm:}
Reconstructing employment measures. Audit the Card and Krueger sample, reproduce descriptive results, compare employment measures, and assess their validity. Submit a pdf report and your R script.

\subsection*{Week 3: Probability as a model of the world.}
\subsubsection*{Class 3.1 (Tuesday 1/19).}
\begin{itemize}
\item How do we think about what it means for something to be random? Why is this mathematical abstraction useful for social science research?
\item Overview of discrete probability.
\item Conditional probability, Bayes Rule.
\item Why understanding conditional probability is essential to understanding social phenomena, and why relying on your intuition can be wrong.
\end{itemize}

\subsubsection*{Readings.}
\begin{itemize}
\item Stark, \textit{SticiGui}, \href{https://www.stat.berkeley.edu/~stark/SticiGui/Text/probabilityAxioms.htm}{Chapter 17: Probability: Axioms and Fundaments}. Read ``The Axioms of Probability'' and ``Conditioning,'' including the subsections ``Conditional Probability,'' ``Independence,'' ``The Multiplication Rule,'' and ``Bayes' Rule.'' Include the worked examples; interactive exercises and videos are optional.
\end{itemize}

\subsubsection*{Class 3.2 (Thursday 1/21).}
\begin{itemize}
\item Using Angrist \& Krueger data to discuss identifying patterns in data, covariance, correlation.
\end{itemize}

\subsubsection*{Readings.}
\begin{itemize}
\item \bibentry{angrist1991does}. Introduction and Figures I--II (p. 983), with their discussion on pp. 982--984.
\item \textbf{Optional:} Stark, \textit{SticiGui}, \href{https://www.stat.berkeley.edu/~stark/SticiGui/Text/association.htm}{Chapter 6: Association}: ``Association of Pairs of Variables'' and ``Post Hoc Ergo Propter Hoc''; \href{https://www.stat.berkeley.edu/~stark/SticiGui/Text/correlation.htm}{Chapter 7: Correlation and Association}: ``The Correlation Coefficient.''
\end{itemize}

\textbf{Homework 3 due Friday 1/22 at 11:59pm:}
Probability, birth cohorts, and schooling. Interpret the published figures, calculate conditional probabilities, and explain how cohort composition can reverse a comparison.

\subsection*{Week 4: Uncertainty and inference.}
\subsubsection*{Class 4.1 (Tuesday 1/26).}
\begin{itemize}
\item Sampling from a larger population; central limit theorem, one- and two-sample t-tests.
\item Confidence intervals and sampling variation.
\end{itemize}

\subsubsection*{Readings.}
\begin{itemize}
\item Stark, \textit{SticiGui}, \href{https://www.stat.berkeley.edu/~stark/SticiGui/Text/standardError.htm}{Chapter 22: Standard Error}. Read the opening definition of standard error and the subsection ``The SE of the Sample Mean of n random Draws from a Box of numbered Tickets,'' including ``The Square-Root Law'' and the comparison of SE with SD.
\item Stark, \textit{SticiGui}, \href{https://www.stat.berkeley.edu/~stark/SticiGui/Text/confidenceIntervals.htm}{Chapter 26: Confidence Intervals}. Read the opening explanation of confidence intervals and coverage, followed by ``Approximate Confidence Intervals for the Population Mean.'' Read the text and worked examples; interactive exercises and videos are optional.
\item \textbf{Optional:} Stark, \textit{SticiGui}, \href{https://www.stat.berkeley.edu/~stark/SticiGui/Text/zTest.htm}{Chapter 30: Approximate Hypothesis Tests: the z Test and the t Test}: ``t Tests'' and ``Hypothesis Tests and Confidence Intervals.''
\end{itemize}

\subsubsection*{Class 4.2 (Thursday 1/28).}
\begin{itemize}
\item Using Butler and Broockman data, randomization inference, introduction to hypothesis testing, confidence intervals, p-values and statistical significance.
\end{itemize}

\subsubsection*{Readings.}
\begin{itemize}
\item \bibentry{butler2011politicians}. Box 1 and design on pp. 466--467, and Table 1 (p. 469).
\item \bibentry{gerber2012field}. Chapter 3: Sampling Distributions, Statistical Inference, and Hypothesis Testing, Sections 3.1--3.5, pp. 51--70.
\item \textbf{Optional:} Stark, \textit{SticiGui}, \href{https://www.stat.berkeley.edu/~stark/SticiGui/Text/testing.htm}{Chapter 27: Hypothesis Testing: Does Chance explain the Results?}: ``Examples of Hypothesis Testing Problems,'' ``Test Statistics and P-values'' and ``Interpreting P-values.''
\end{itemize}

\textbf{Homework 4 due Friday 1/29 at 11:59pm:}
Assignment mechanisms and uncertainty. Construct exact null distributions, compare assignment procedures, and interpret published evidence and confidence intervals.

\subsection*{Week 5: Midterm.}
\subsubsection*{Class 5.1 (Tuesday 2/2).}
\begin{itemize}
\item Midterm synthesis and review.
\end{itemize}

\subsubsection*{Readings.}
\begin{itemize}
\item Rework Homeworks 1--4 and the midterm practice distributed on Canvas. No new readings.
\end{itemize}

\subsubsection*{Class 5.2 (Thursday 2/4).}
\begin{itemize}
\item In-class midterm.
\end{itemize}

\subsubsection*{Readings.}
\begin{itemize}
\item No readings.
\end{itemize}

\subsection*{Week 6: Bivariate regression: best linear predictor.}
\subsubsection*{Class 6.1 (Tuesday 2/9).}
\begin{itemize}
\item With Butler and Broockman data, regression as the best linear predictor, least squares visually demonstrated.
\item Interpreting regression coefficients.
\end{itemize}

\subsubsection*{Readings.}
\begin{itemize}
\item \bibentry{buenodemesquita2021thinking}. Chapter 5: Regression for Describing and Forecasting.
\end{itemize}

\subsubsection*{Class 6.2 (Thursday 2/11).}
\begin{itemize}
\item Inference for linear regression.
\item Linear regression with experiments and dummy variables; relationship to t-tests.
\end{itemize}

\subsubsection*{Readings.}
\begin{itemize}
\item \bibentry{pager2003mark}. Design on pp. 946--947 and Figure 6 (p. 958).
\item \textbf{Optional:} Stark, \textit{SticiGui}, \href{https://www.stat.berkeley.edu/~stark/SticiGui/Text/zTest.htm}{Chapter 30: Approximate Hypothesis Tests: the z Test and the t Test}: ``t Tests'' and ``Hypothesis Tests and Confidence Intervals.''
\end{itemize}

\textbf{Homework 5 due Friday 2/12 at 11:59pm:}
Regression, evidence, and prediction. Derive and verify binary-predictor OLS results with Card and Krueger data, interpret robust uncertainty, and apply the identity to Pager. Submit a pdf report and your R script.

\subsection*{Week 7: Multivariate regression: model building.}
\subsubsection*{Class 7.1 (Tuesday 2/16).}
\begin{itemize}
\item Regression as linearization of conditional expectation function.
\item With Banerjee et al. data, interpreting regression coefficients (and standard errors) in multivariate regression.
\end{itemize}

\subsubsection*{Readings.}
\begin{itemize}
\item \bibentry{banerjee2015miracle}. Introduction, interpretation and regression discussion on pp. 35--36, and Table 2 (p. 37).
\end{itemize}

\subsubsection*{Class 7.2 (Thursday 2/18).}
\begin{itemize}
\item When can we apply causal interpretations to regression coefficients from observational data?
\end{itemize}

\subsubsection*{Readings.}
\begin{itemize}
\item \bibentry{cardkrueger1994minimum}.
\begin{itemize}
\item Return to Section III.A and Table 3 (p. 780) to examine their approach to analysis.
\end{itemize}
\item \textbf{Optional:} Stark, \textit{SticiGui}, \href{https://www.stat.berkeley.edu/~stark/SticiGui/Text/experiments.htm}{Chapter 28: Does Treatment Have an Effect?}: ``The Method of Comparison,'' ``Confounding'' and ``Experiments and Observational Studies.''
\end{itemize}

\textbf{Homework 6 due Friday 2/19 at 11:59pm:}
Controls and competing causal designs. Compare Banerjee et al.'s randomized design with Card and Krueger's trend comparison, including adjustment and sensitivity.

\subsection*{Week 8: Statistical power and pre-registration.}
\subsubsection*{Class 8.1 (Tuesday 2/23).}
\begin{itemize}
\item Statistical power, effect sizes, sample size.
\end{itemize}

\subsubsection*{Readings.}
\begin{itemize}
\item \bibentry{coppock2013power}. \url{https://methods.egap.org/guides/assessing-designs/power_en.html}
\item \bibentry{athey2017econometrics}. Read Section 7.1, ``Power calculations,'' pp. 102--104. Focus on planning inputs and the equal-arm sample-size relationship; no derivation is required.
\item \textbf{Optional:} Stark, \textit{SticiGui}, \href{https://www.stat.berkeley.edu/~stark/SticiGui/Text/testing.htm}{Chapter 27}: ``Significance Level and Power'' and ``Statistical Significance and Practical Importance.''
\end{itemize}

\subsubsection*{Class 8.2 (Thursday 2/25).}
\begin{itemize}
\item Declaring hypotheses in advance / pre-analysis planning.
\end{itemize}

\subsubsection*{Readings.}
\begin{itemize}
\item \bibentry{chengradypreanalysis}. \url{https://methods.egap.org/guides/planning/pap_en.html}
\item \bibentry{noy2023experimental}. Abstract and research design.
\item \bibentry{noy2023planned}. \href{https://www.socialscienceregistry.org/versions/170369/docs/version/document}{Registered plan.} Read ``Outcome variables (at the task level)'' and ``Productivity Main Effects.'' Focus on the planned outcomes and comparisons; detailed regression syntax is not required.
\item \textbf{Optional:} Stark, \textit{SticiGui}, \href{https://www.stat.berkeley.edu/~stark/SticiGui/Text/testing.htm}{Chapter 27}: ``Multiplicity and Data Mining.''
\end{itemize}

\textbf{Homework 7 due Friday 2/26 at 11:59pm:}
Power and preanalysis planning. Audit Noy and Zhang's registered plan, assess recruitment and allocation choices, and write a feasible follow-up plan.

\subsection*{Week 9: Review and cumulative final.}
\textbf{Homework 8 due Tuesday 3/2 at 9:00am:}
Compare evidence about employment. Evaluate the magnitudes, uncertainty, and identifying arguments in Pager and Card and Krueger, then propose a follow-up study in a policy memo.

\subsubsection*{Class 9.1 (Tuesday 3/2).}
\begin{itemize}
\item Cumulative synthesis and final review.
\end{itemize}

\subsubsection*{Readings.}
\begin{itemize}
\item Rework Homework 8 and the final practice distributed on Canvas. No new readings.
\end{itemize}

\subsubsection*{Class 9.2 (Thursday 3/4).}
\begin{itemize}
\item In-class cumulative final.
\end{itemize}

\subsubsection*{Readings.}
\begin{itemize}
\item No readings.
\end{itemize}


\end{document}
